% chapters/17/section-01.tex
\section{Peta Kompetensi dari CPL sampai Sub-CPMK}
\label{sec:peta-kompetensi}

\subsection{Satu Rantai, Empat Tingkat}

Seluruh mata kuliah ini disusun sebagai satu rantai. Di puncaknya ada dua
capaian pembelajaran lulusan yang dibebankan kepada mata kuliah ini. Di
bawahnya ada empat capaian pembelajaran mata kuliah, lalu tiga belas
Sub-CPMK, dan di dasar rantai itu tujuh belas bab yang Anda pelajari.

Rantai itu bukan hiasan administratif. Setiap butir pada rantai menjawab satu
pertanyaan, dan pertanyaannya berbeda pada setiap tingkat.

\begin{description}
  \item[CPL] Apa yang harus dimiliki seorang lulusan program studi ini?
  \item[CPMK] Bagian mana dari kemampuan itu yang dibangun oleh mata kuliah
        ini?
  \item[Sub-CPMK] Bagian mana dari CPMK itu yang dibangun oleh satu pokok
        bahasan?
  \item[Indikator] Perilaku apa yang dapat diamati untuk menyatakan bahwa
        pokok bahasan itu sudah dikuasai?
\end{description}

Rantai itu diperlihatkan pada Gambar~\ref{fig:peta-kompetensi}. Angka di
bawah setiap CPMK adalah bobotnya terhadap nilai akhir, dan angka di bawah
setiap CPL adalah jumlah bobot CPMK penyusunnya.

\begin{figure}[htbp]
\centering
% chapters/17/figures/peta-kompetensi.tikz
% Peta capaian: dari CPL, turun ke CPMK, lalu ke kelompok bab yang
% menuntaskannya. Hanya keterkaitan utamanya yang digambar; setiap CPMK
% juga didukung oleh CPL yang lain.
\begin{tikzpicture}[
  cpl/.style={
    rectangle, rounded corners=5pt, draw=ubPrimary, thick, fill=ubPrimary,
    text=white, align=center, text width=30mm, inner sep=4pt,
    minimum height=11mm, font=\footnotesize\bfseries
  },
  cpmk/.style={
    rectangle, rounded corners=5pt, draw=ubPrimary, thick, fill=ubAlt,
    align=center, text width=28mm, inner sep=4pt, minimum height=11mm,
    font=\footnotesize
  },
  bab/.style={
    rectangle, rounded corners=5pt, draw=ubAccent, thick, fill=white,
    align=center, text width=28mm, inner sep=4pt, minimum height=9mm,
    font=\footnotesize
  },
  panah/.style={-{Stealth[length=2.2mm]}, thick, draw=ubAccent},
  bobot/.style={font=\scriptsize, ubPrimary}
]

\node[cpl] (cpl3) at (1.8,2.6) {\CPL{03}\\{\scriptsize\itshape 57\%}};
\node[cpl] (cpl4) at (8.4,2.6) {\CPL{04}\\{\scriptsize\itshape 43\%}};

\node[cpmk] (c1) at (0,1.0) {\CPMK{1}\\{\scriptsize\itshape 25\%}};
\node[cpmk] (c2) at (3.6,1.0) {\CPMK{2}\\{\scriptsize\itshape 32\%}};
\node[cpmk] (c3) at (6.6,1.0) {\CPMK{3}\\{\scriptsize\itshape 22\%}};
\node[cpmk] (c4) at (10.2,1.0) {\CPMK{4}\\{\scriptsize\itshape 21\%}};

\node[bab] (b1) at (0,-0.7) {Bab III--VII\\{\scriptsize\itshape Sub-CPMK 1.1--1.4}};
\node[bab] (b2) at (3.6,-0.7) {Bab VIII--XI\\{\scriptsize\itshape Sub-CPMK 2.1--2.4}};
\node[bab] (b3) at (6.6,-0.7) {Bab XII--XIII\\{\scriptsize\itshape Sub-CPMK 3.1--3.2}};
\node[bab] (b4) at (10.2,-0.7) {Bab XIV--XVI\\{\scriptsize\itshape Sub-CPMK 4.1--4.3}};

\draw[panah] (cpl3) -- (c1);
\draw[panah] (cpl3) -- (c2);
\draw[panah] (cpl4) -- (c3);
\draw[panah] (cpl4) -- (c4);
\draw[panah] (c1) -- (b1);
\draw[panah] (c2) -- (b2);
\draw[panah] (c3) -- (b3);
\draw[panah] (c4) -- (b4);
\end{tikzpicture}

\caption{Peta capaian dari CPL sampai kelompok bab yang menuntaskannya}
\label{fig:peta-kompetensi}
\end{figure}

Gambar itu hanya memperlihatkan keterkaitan utamanya. Pada RPS, setiap CPMK
juga didukung oleh CPL yang lain: \CPMK{1} dan \CPMK{2} bersandar pada
\CPL{03} sebagai utama dan \CPL{04} sebagai pendukung, sedangkan \CPMK{3} dan
\CPMK{4} berlaku sebaliknya.

\subsection{Tiga Belas Sub-CPMK dan Bab yang Menuntaskannya}

Tabel~\ref{tab:peta-subcpmk} memuat seluruh Sub-CPMK beserta bab dan instrumen
yang mengukurnya. Tabel ini adalah ringkasan dari seluruh buku, dan sebaiknya
Anda buka kembali sebelum menempuh Ujian Akhir Semester.

\begin{table}[htbp]
\centering
\caption{Seluruh Sub-CPMK, bab yang menuntaskannya, dan instrumen penilaiannya}
\label{tab:peta-subcpmk}
{\setlength{\tabcolsep}{4pt}\footnotesize%
\begin{tabularx}{\textwidth}{@{}C{1.4cm} C{1.8cm} L{3.1cm} C{1.5cm} Y@{}}
\toprule
\headrow \textbf{Sub-CPMK} & \textbf{CPMK} & \textbf{Pokok bahasan} & \textbf{Bab} & \textbf{Instrumen} \\
\midrule
1.1 & 1 & \istilah{Game loop}, \istilah{event handling}, \istilah{clock} & IV & Praktikum 02, Kuis 1, UTS \\
\addlinespace
1.2 & 1 & Gerak objek, vektor, masukan, \istilah{delta time} & V & Praktikum 03, Kuis 1, UTS \\
\addlinespace
1.3 & 1 & Fisika sederhana & VI & Praktikum 04, UTS \\
\addlinespace
1.4 & 1 & \istilah{Sprite}, kelompok, deteksi tabrakan & VII & Praktikum 05, UTS \\
\addlinespace
2.1 & 2 & AI NPC dan perilaku pengejar & VIII & Praktikum 06, Kuis 2 \\
\addlinespace
2.2 & 2 & Mesin keadaan untuk NPC & IX & Praktikum 07, Kuis 2 \\
\addlinespace
2.3 & 2 & Pencarian jalur pada petak: BFS, Dijkstra, A* & X & Praktikum 08, Kuis 2 \\
\addlinespace
2.4 & 2 & Integrasi mesin keadaan dan A* & XI & Mini Proyek AI \\
\addlinespace
3.1 & 3 & Unity: \istilah{project}, \istilah{scene}, \istilah{GameObject} & XII & Praktikum 09 \\
\addlinespace
3.2 & 3 & Level, fisika, dan kendali pemain & XIII & Praktikum 10, Proyek Akhir \\
\addlinespace
4.1 & 4 & Aset visual, animasi, dan antarmuka & XIV & Praktikum 11, Proyek Akhir \\
\addlinespace
4.2 & 4 & Audio dan pemicu interaksi & XV & Praktikum 12, Proyek Akhir \\
\addlinespace
4.3 & 4 & Evaluasi prototipe dan presentasi & XVI & Proyek Akhir, UAS \\
\bottomrule
\end{tabularx}}
\end{table}

\begin{catatabox}
Setiap Sub-CPMK memiliki sekurang-kurangnya dua instrumen.
Itu disengaja. Satu instrumen dapat keliru mengukur, sedangkan dua instrumen
yang berbeda cara mengukurnya jarang keliru bersama-sama. Apabila nilai
praktikum Anda tinggi tetapi nilai ujian Anda rendah, yang perlu diperiksa
bukan penilaiannya, melainkan apakah pekerjaan praktikum itu benar-benar Anda
kerjakan sendiri.
\end{catatabox}
